% SOB-ID FSM (TikZ)
% Usage in LaTeX:
%   \usepackage{tikz}
%   \usetikzlibrary{arrows.meta,positioning}
%   ...
%   \begin{figure}[t]
%   \centering
%   % SOB-ID FSM (TikZ)
% Usage in LaTeX:
%   \usepackage{tikz}
%   \usetikzlibrary{arrows.meta,positioning}
%   ...
%   \begin{figure}[t]
%   \centering
%   % SOB-ID FSM (TikZ)
% Usage in LaTeX:
%   \usepackage{tikz}
%   \usetikzlibrary{arrows.meta,positioning}
%   ...
%   \begin{figure}[t]
%   \centering
%   % SOB-ID FSM (TikZ)
% Usage in LaTeX:
%   \usepackage{tikz}
%   \usetikzlibrary{arrows.meta,positioning}
%   ...
%   \begin{figure}[t]
%   \centering
%   \input{sob_id_fsm.tex}
%   \caption{...}\label{fig:sobid-fsm}
%   \end{figure}

\begin{tikzpicture}[
  >=Stealth,
  node distance=18mm and 18mm,
  every node/.style={draw, rounded corners, align=left, font=\small, inner sep=3pt},
  arrow/.style={->, thick}
]

\node (s0) {s0: login\_page};
\node (s1) [right=of s0] {s1: auth\_attempted\\captcha\_challenge};
\node (s2) [right=of s1, yshift=15mm] {s2: mfa\_sent};
\node (s3) [right=of s1, yshift=-15mm] {s3: session\_active\\device\_trusted};
\node (s4) [below=of s1] {s4: authentication\_failed};

\node (s5) [right=of s2] {s5: geolocation\_mismatch\\ip\_blacklisted\\anomaly\_score\_high};
\node (s6) [right=of s3] {s6: threshold\_breached\\transaction\_flagged\\alert\_triggered\\anomaly\_score\_high};
\node (s11) [right=of s6, yshift=15mm] {s11: malware\_detected\\alert\_triggered};

\node (s7) [right=of s6, yshift=-15mm] {s7: manual\_review\_requested};
\node (s8) [right=of s7] {s8: account\_locked\\alert\_triggered};
\node (s10) [below=of s7] {s10: password\_reset\_required};
\node (s9) [below=of s8] {s9: session\_terminated};

% Start arrow (initial state is s0)
\node (start) [left=10mm of s0, draw=none] {};
\draw[arrow] (start) -- (s0);

% Edges (exactly match the NuSMV model)
\draw[arrow] (s0) -- (s1);

\draw[arrow] (s1) -- (s2);
\draw[arrow] (s1) -- (s3);
\draw[arrow] (s1) -- (s4);

\draw[arrow] (s2) -- (s3);
\draw[arrow] (s2) -- (s4);
\draw[arrow] (s2) -- (s5);

\draw[arrow] (s3) -- (s6);
\draw[arrow] (s3) -- (s11);
\draw[arrow] (s3) -- (s9);
\draw[arrow] (s3) edge[loop below] (s3);

\draw[arrow] (s4) -- (s1);
\draw[arrow] (s4) -- (s8);

\draw[arrow] (s5) -- (s6);
\draw[arrow] (s5) -- (s7);
\draw[arrow] (s5) -- (s3);

\draw[arrow] (s6) -- (s7);
\draw[arrow] (s6) -- (s8);

\draw[arrow] (s7) -- (s3);
\draw[arrow] (s7) -- (s8);
\draw[arrow] (s7) -- (s10);

\draw[arrow] (s8) -- (s9);
\draw[arrow] (s9) -- (s0);

\draw[arrow] (s10) -- (s1);
\draw[arrow] (s10) -- (s9);

\draw[arrow] (s11) -- (s7);
\draw[arrow] (s11) -- (s8);

\end{tikzpicture}

%   \caption{...}\label{fig:sobid-fsm}
%   \end{figure}

\begin{tikzpicture}[
  >=Stealth,
  node distance=18mm and 18mm,
  every node/.style={draw, rounded corners, align=left, font=\small, inner sep=3pt},
  arrow/.style={->, thick}
]

\node (s0) {s0: login\_page};
\node (s1) [right=of s0] {s1: auth\_attempted\\captcha\_challenge};
\node (s2) [right=of s1, yshift=15mm] {s2: mfa\_sent};
\node (s3) [right=of s1, yshift=-15mm] {s3: session\_active\\device\_trusted};
\node (s4) [below=of s1] {s4: authentication\_failed};

\node (s5) [right=of s2] {s5: geolocation\_mismatch\\ip\_blacklisted\\anomaly\_score\_high};
\node (s6) [right=of s3] {s6: threshold\_breached\\transaction\_flagged\\alert\_triggered\\anomaly\_score\_high};
\node (s11) [right=of s6, yshift=15mm] {s11: malware\_detected\\alert\_triggered};

\node (s7) [right=of s6, yshift=-15mm] {s7: manual\_review\_requested};
\node (s8) [right=of s7] {s8: account\_locked\\alert\_triggered};
\node (s10) [below=of s7] {s10: password\_reset\_required};
\node (s9) [below=of s8] {s9: session\_terminated};

% Start arrow (initial state is s0)
\node (start) [left=10mm of s0, draw=none] {};
\draw[arrow] (start) -- (s0);

% Edges (exactly match the NuSMV model)
\draw[arrow] (s0) -- (s1);

\draw[arrow] (s1) -- (s2);
\draw[arrow] (s1) -- (s3);
\draw[arrow] (s1) -- (s4);

\draw[arrow] (s2) -- (s3);
\draw[arrow] (s2) -- (s4);
\draw[arrow] (s2) -- (s5);

\draw[arrow] (s3) -- (s6);
\draw[arrow] (s3) -- (s11);
\draw[arrow] (s3) -- (s9);
\draw[arrow] (s3) edge[loop below] (s3);

\draw[arrow] (s4) -- (s1);
\draw[arrow] (s4) -- (s8);

\draw[arrow] (s5) -- (s6);
\draw[arrow] (s5) -- (s7);
\draw[arrow] (s5) -- (s3);

\draw[arrow] (s6) -- (s7);
\draw[arrow] (s6) -- (s8);

\draw[arrow] (s7) -- (s3);
\draw[arrow] (s7) -- (s8);
\draw[arrow] (s7) -- (s10);

\draw[arrow] (s8) -- (s9);
\draw[arrow] (s9) -- (s0);

\draw[arrow] (s10) -- (s1);
\draw[arrow] (s10) -- (s9);

\draw[arrow] (s11) -- (s7);
\draw[arrow] (s11) -- (s8);

\end{tikzpicture}

%   \caption{...}\label{fig:sobid-fsm}
%   \end{figure}

\begin{tikzpicture}[
  >=Stealth,
  node distance=18mm and 18mm,
  every node/.style={draw, rounded corners, align=left, font=\small, inner sep=3pt},
  arrow/.style={->, thick}
]

\node (s0) {s0: login\_page};
\node (s1) [right=of s0] {s1: auth\_attempted\\captcha\_challenge};
\node (s2) [right=of s1, yshift=15mm] {s2: mfa\_sent};
\node (s3) [right=of s1, yshift=-15mm] {s3: session\_active\\device\_trusted};
\node (s4) [below=of s1] {s4: authentication\_failed};

\node (s5) [right=of s2] {s5: geolocation\_mismatch\\ip\_blacklisted\\anomaly\_score\_high};
\node (s6) [right=of s3] {s6: threshold\_breached\\transaction\_flagged\\alert\_triggered\\anomaly\_score\_high};
\node (s11) [right=of s6, yshift=15mm] {s11: malware\_detected\\alert\_triggered};

\node (s7) [right=of s6, yshift=-15mm] {s7: manual\_review\_requested};
\node (s8) [right=of s7] {s8: account\_locked\\alert\_triggered};
\node (s10) [below=of s7] {s10: password\_reset\_required};
\node (s9) [below=of s8] {s9: session\_terminated};

% Start arrow (initial state is s0)
\node (start) [left=10mm of s0, draw=none] {};
\draw[arrow] (start) -- (s0);

% Edges (exactly match the NuSMV model)
\draw[arrow] (s0) -- (s1);

\draw[arrow] (s1) -- (s2);
\draw[arrow] (s1) -- (s3);
\draw[arrow] (s1) -- (s4);

\draw[arrow] (s2) -- (s3);
\draw[arrow] (s2) -- (s4);
\draw[arrow] (s2) -- (s5);

\draw[arrow] (s3) -- (s6);
\draw[arrow] (s3) -- (s11);
\draw[arrow] (s3) -- (s9);
\draw[arrow] (s3) edge[loop below] (s3);

\draw[arrow] (s4) -- (s1);
\draw[arrow] (s4) -- (s8);

\draw[arrow] (s5) -- (s6);
\draw[arrow] (s5) -- (s7);
\draw[arrow] (s5) -- (s3);

\draw[arrow] (s6) -- (s7);
\draw[arrow] (s6) -- (s8);

\draw[arrow] (s7) -- (s3);
\draw[arrow] (s7) -- (s8);
\draw[arrow] (s7) -- (s10);

\draw[arrow] (s8) -- (s9);
\draw[arrow] (s9) -- (s0);

\draw[arrow] (s10) -- (s1);
\draw[arrow] (s10) -- (s9);

\draw[arrow] (s11) -- (s7);
\draw[arrow] (s11) -- (s8);

\end{tikzpicture}

%   \caption{...}\label{fig:sobid-fsm}
%   \end{figure}

\begin{tikzpicture}[
  >=Stealth,
  node distance=18mm and 18mm,
  every node/.style={draw, rounded corners, align=left, font=\small, inner sep=3pt},
  arrow/.style={->, thick}
]

\node (s0) {s0: login\_page};
\node (s1) [right=of s0] {s1: auth\_attempted\\captcha\_challenge};
\node (s2) [right=of s1, yshift=15mm] {s2: mfa\_sent};
\node (s3) [right=of s1, yshift=-15mm] {s3: session\_active\\device\_trusted};
\node (s4) [below=of s1] {s4: authentication\_failed};

\node (s5) [right=of s2] {s5: geolocation\_mismatch\\ip\_blacklisted\\anomaly\_score\_high};
\node (s6) [right=of s3] {s6: threshold\_breached\\transaction\_flagged\\alert\_triggered\\anomaly\_score\_high};
\node (s11) [right=of s6, yshift=15mm] {s11: malware\_detected\\alert\_triggered};

\node (s7) [right=of s6, yshift=-15mm] {s7: manual\_review\_requested};
\node (s8) [right=of s7] {s8: account\_locked\\alert\_triggered};
\node (s10) [below=of s7] {s10: password\_reset\_required};
\node (s9) [below=of s8] {s9: session\_terminated};

% Start arrow (initial state is s0)
\node (start) [left=10mm of s0, draw=none] {};
\draw[arrow] (start) -- (s0);

% Edges (exactly match the NuSMV model)
\draw[arrow] (s0) -- (s1);

\draw[arrow] (s1) -- (s2);
\draw[arrow] (s1) -- (s3);
\draw[arrow] (s1) -- (s4);

\draw[arrow] (s2) -- (s3);
\draw[arrow] (s2) -- (s4);
\draw[arrow] (s2) -- (s5);

\draw[arrow] (s3) -- (s6);
\draw[arrow] (s3) -- (s11);
\draw[arrow] (s3) -- (s9);
\draw[arrow] (s3) edge[loop below] (s3);

\draw[arrow] (s4) -- (s1);
\draw[arrow] (s4) -- (s8);

\draw[arrow] (s5) -- (s6);
\draw[arrow] (s5) -- (s7);
\draw[arrow] (s5) -- (s3);

\draw[arrow] (s6) -- (s7);
\draw[arrow] (s6) -- (s8);

\draw[arrow] (s7) -- (s3);
\draw[arrow] (s7) -- (s8);
\draw[arrow] (s7) -- (s10);

\draw[arrow] (s8) -- (s9);
\draw[arrow] (s9) -- (s0);

\draw[arrow] (s10) -- (s1);
\draw[arrow] (s10) -- (s9);

\draw[arrow] (s11) -- (s7);
\draw[arrow] (s11) -- (s8);

\end{tikzpicture}
